\documentclass[11pt,a4paper]{article}
\usepackage[T2A,T1]{fontenc}
\usepackage[utf8]{inputenc}
\usepackage[english,russian]{babel}
\usepackage[margin=2.2cm]{geometry}
\usepackage{amsmath,amssymb}
\usepackage{xcolor}
\usepackage{titlesec}

\titleformat{\section}{\normalfont\large\bfseries\color{blue!50!black}}
            {\thesection.}{0.5em}{}
\titlespacing*{\section}{0pt}{1.0em}{0.3em}

\setlength{\parskip}{0.35em}
\setlength{\parindent}{0pt}

\title{\vspace{-2em}\textbf{0--k BFS -- Текст доклада}\\
\large Сопровождение к \texttt{presentation\_ru.pdf}}
\author{Артёмов Макар, Марк Шкут \\ Сколтех, Алгоритмы 2026}
\date{21 мая 2026}

\begin{document}
\maketitle
\vspace{-1em}

Ниже -- буквальный текст того, что нужно говорить на каждом слайде.
Один раздел -- один слайд; зачитываем абзац и переключаем слайд, дойдя
до пунктирной линии.

\section{Титульный слайд}
``Добрый день. Меня зовут Макар, со мной Марк. Мы реализовали на
C++17 семейство алгоритмов 0--k BFS для задачи кратчайших путей и
сравнили их с Дейкстрой. Заголовочный результат: на графе с 900 миллионами
вершин наш алгоритм отрабатывает за две минуты, а Дейкстра за десять. Я
расскажу алгоритм, покажу его в работе и пройдусь по замерам. Минут
пятнадцать плюс вопросы.''
\hrulefill

\section{Задача}
``Задача -- кратчайшие пути из одной вершины во взвешенном графе. Особенность
в том, что веса -- неотрицательные целые числа, ограниченные сверху
небольшим известным $k$. Это ограничение не искусственное: оно
постоянно встречается в планировании пути на сетках, в слоистых графах,
при сведениях к 2-SAT, в транспортных сетях с несколькими уровнями
стоимости. Конкретная цель -- довести размер входа до миллиарда вершин
на одной рабочей станции, насколько позволяет железо.''
\hrulefill

\section{Базовый алгоритм: Дейкстра}
``Классический ответ -- Дейкстра. Реализация учебная: \texttt{std::priority\_queue}
из пар (расстояние, вершина) и ленивое decrease-key -- при каждом
улучшении расстояния кладём свежую запись, а при поп пропускаем
устаревшие. Сложность $O((V+E) \log V)$. Работает для любых неотрицательных
вещественных весов. Множитель $\log V$ платится при каждой успешной
релаксации -- это и есть цена кучи. Вопрос: можем ли мы избавиться от
него, если веса -- маленькие целые числа?''
\hrulefill

\section{Идея: 0--1 BFS}
``Простейший случай -- веса 0 или 1. Заменяем кучу на deque. Ребро веса
0 кладём в начало, ребро веса 1 -- в конец. Deque всё время упорядочен
по расстоянию и содержит максимум два различных значения -- текущий
слой спереди, следующий сзади. Линейное время, линейная память. Тонкость
по сравнению с обычным BFS: вершина может быть открыта несколько раз
до того, как её расстояние станет окончательным, поэтому проверяем не
``посещена/нет'', а $dist[v] + w < dist[u]$.''
\hrulefill

\section{Обобщение: 0--k BFS}
``Теперь обобщаем на веса до $k$. Наблюдение: релаксация двигает вершину
максимум на $k$ шагов вперёд в порядке расстояний. Значит, в любой момент
живых слоёв расстояний не больше $k+1$. Заводим $k+1$ FIFO-очередь и
индексируем их по расстоянию по модулю $k+1$ -- это кольцевой буфер
бакетов. Указатель \texttt{cur} обходит их. Релаксация кладёт цель в
бакет нового расстояния. Когда активный бакет опустеет -- \texttt{cur}
сдвигается вперёд. Устаревшие записи от прошлых релаксаций остаются в
своих старых бакетах и пропускаются при поп. Сложность $O(kV+E)$, память
$O(k+E)$ на очереди плюс $O(V)$ на массив расстояний.''
\hrulefill

\section{Как это выглядит}
``Вот оба алгоритма на одной и той же сетке -- 200 на 200, $k=8$. Они
выдают одинаковое поле расстояний -- волна расходится радиально от
центра, форму задаёт случайный рельеф. В заголовках -- реально замеренное
время. Левая панель, 0--k BFS, закончила за четыре миллисекунды. Правая
панель, Дейкстра -- за тринадцать с третью миллисекунд -- в три с лишним
раза дольше. Та же корректность, существенно меньше времени. Сама анимация
есть в репозитории, здесь -- один кадр.''
\hrulefill

\section{Под капотом}
``Это педагогический центр. На крошечной сетке -- десять на десять, $k=3$ --
мы инструментировали алгоритм и выгружаем каждое событие. Справа -- четыре
реальные очереди $Q[0]$\,--\,$Q[3]$ с их текущим содержимым; числа в клетках --
идентификаторы вершин. Слева сетка; цвет клетки показывает, в каком бакете
сейчас лежит эта вершина. Красная стрелка \texttt{cur} указывает на активный
бакет. Видно, как релаксация кладёт вершину в бакет, выбранный новым
расстоянием по модулю четыре, и как активный бакет вычерпывается в порядке
FIFO до того, как \texttt{cur} сдвинется. Это весь секрет алгоритма.''
\hrulefill

\section{Дизайн под V $\to 10^9$}
``Теперь про инженерию. Чтобы дойти до миллиарда вершин, нельзя хранить
рёбра явно: CSR с четырьмя рёбрами на вершину и 32-битными
идентификаторами стоил бы уже двадцать гигабайт. Поэтому мы сделали два
бэкенда. Универсальный \texttt{GraphCSR} для обычных графов и неявный
\texttt{GridGraph}, который вычисляет четырёх соседей по запросу, без
хранения рёбер. Второе -- сделали тип расстояния \texttt{uint32\_t}
вместо \texttt{uint64\_t}, что уменьшает доминирующий массив вдвое.
Компромисс допустим, пока $k \cdot V$ помещается в 32 бита, а для наших
параметров это так; есть и 64-битный fallback.''
\hrulefill

\section{Результаты}
``Вот цифры -- всё на сетке, в один поток. На миллионе вершин уже
получаем ускорение в 4.6 раза. На десяти миллионах -- 5.6. На ста
миллионах при $k=1$ deque-вариант 0--1 BFS даёт одиннадцатикратное
ускорение -- это пик. Следующая строка интересна: на половине миллиарда
вершин оба алгоритма становятся memory-bandwidth bound и разрыв
схлопывается до единицы. Внизу заголовочный результат: на 900 миллионах
вершин 0--1 BFS делает работу за 121 секунду, Дейкстра -- за 620.
Все контрольные суммы расстояний совпадают между алгоритмами --
19 различных инстансов в validation sweep, ноль расхождений.''
\hrulefill

\section{Масштабирование}
``Те же данные, log-log график. Обе оси логарифмические. Дейкстра --
верхняя линия, варианты BFS идут параллельно ниже. Постоянный вертикальный
зазор в логарифмической шкале означает постоянное мультипликативное
ускорение. Справа, при $10^8$ и выше, линии начинают сходиться -- мы
входим в memory-bound режим. Самая нижняя линия всегда -- deque-вариант
0--1 BFS, он чисто впереди.''
\hrulefill

\section{Ограничения}
``Честно про границы. Первое: это выигрыш по константе, не асимптотический.
Как только $k$ сравним с $\log V$, Дейкстра не хуже. Второе: при $V$
больше полумиллиарда всё упирается в память, и накладные расходы очередей
начинают мешать. Третье: содержать $k+1$ отдельный FIFO-аллокатор
дороже для кэша, чем один deque -- именно поэтому при $V \approx 10^9$
наш 0--1 BFS всё ещё выигрывает, а общий 0--k BFS на этом масштабе уже
нет. Четвёртое: это не замена Дейкстре в общем случае; для вещественных
весов или неизвестного диапазона -- использовать Дейкстру.''
\hrulefill

\section{Дальнейшие шаги}
``Четыре пункта в плане. Первый: настоящая индексная $d$-арная куча с
честным decrease-key для baseline -- чтобы сравнение было справедливее.
Второй: кусочный кольцевой FIFO для 0--k BFS, чтобы закрыть разрыв с
0--1 BFS по кэшу. Третий: выйти из зоны синтетических графов и загрузить
реальную дорожную сеть -- DIMACS USA с весами, квантованными до $k=8$.
Четвёртый, стретч-цель: вариант $\Delta$-stepping для параллельного
контекста.''
\hrulefill

\section{Спасибо}
``Всё это -- код, LaTeX-отчёт, анимации, воспроизводимые скрипты сборки --
лежит в GitHub-репозитории по ссылке. Готовы к вопросам.''

\textbf{Вероятные вопросы:}
\begin{itemize}
\item \emph{Почему 0--1 BFS быстрее 0--k BFS при $k=1$?} -- один
\texttt{std::deque} против двух \texttt{std::queue}: накладные расходы
аллокатора и кэша.
\item \emph{Как проверяли корректность на запуске с миллиардом вершин?} --
суммы расстояний совпадают с Дейкстрой побайтово; сумма фиксирует массив,
поэлементное сравнение не требуется.
\item \emph{А параллелизм?} -- весь Этап 1 в один поток;
$\Delta$-stepping -- в плане.
\end{itemize}

\end{document}
